\section{Capacity study: design}
\label{sec:design}

The study asks how fast exact recall collapses as stored patterns grow alike, which rules hold up, and whether the characters fail from density, shared structure, or load. A recall is \emph{exact} when the final state equals the query's own target, so one wrong pixel is a failure.

\subsection{Patterns}

Every network has $N = 144$ neurons. Two families of random $\pm 1$ patterns have controlled overlap. In the \emph{prototype} family, each pattern flips each bit of a shared random prototype independently with probability $(1 - \sqrt{\rho})/2$, so the expected overlap between two patterns is $\rho$ and the expected density is 0.5 (a set's density varies around it); $\rho = 0, 0.1, \dots, 0.9$. In the \emph{biased} family, bits are independent and on with probability 0.65 or 0.80, which gives expected overlaps $(2d-1)^2$ of 0.09 and 0.36 (measured at 40 patterns: \N{po:biased:0.65:4} and \N{po:biased:0.8}). The load $P$ takes the values 2 to 9, then 12, 15, 18, 24, 30, and 40. Each condition has 50 pattern sets, and the network at load $P$ stores the first $P$ patterns of a set of 40. Queries flip 14 or 29 of the 144 bits (labeled 10\% and 20\% throughout; the test file flips \N{d:test_flips}). The classical rules recall with asynchronous updates in random order, for up to 50 sweeps; the modern network and nearest-pattern lookup work as described below. Every rule sees the same patterns and queries, and the classical rules use the same update orders.

\subsection{Rules}

\emph{Hebbian} is the outer-product rule with the standard update, and \emph{Hebbian, 0/1 update} uses the same weights with the assignment's update~\eqref{eq:bias}. \emph{Storkey}'s rule adds each pattern as $W_{ij} \leftarrow W_{ij} + \frac{1}{N}(\xi_i\xi_j - \xi_i h_{ji} - h_{ij}\xi_j)$, with the local field $h_{ij} = \sum_{k \neq i,j} W_{ik}\xi_k$ \cite{storkey1997}. It is incremental, so the stored order matters. The \emph{pseudoinverse} rule is $W = \Xi(\Xi^\top \Xi)^{+}\Xi^\top$ with a zero diagonal \cite{personnaz1985,kanter1987}, where the columns of $\Xi$ are the stored patterns. The \emph{covariance} rule \cite{amit1987} uses $W = \sum_\mu (\xi^\mu - a)(\xi^\mu - a)^\top$, where $a$ is the mean of all stored $\pm 1$ bits ($a = 2d - 1$ at density $d$), computes the field on the centered state, and thresholds near the midpoint between the expected fields of a $+1$ and a $-1$ bit; at $a = 0$, it is the Hebbian rule. The \emph{modern} network \cite{ramsauer2021} on unit vectors updates $x \leftarrow \Xi\,\mathrm{softmax}(\beta\,\Xi^\top x)$ and renormalizes, up to 10 times or until $x$ stops changing, and outputs the sign; the study uses $\beta = 4, 16, 64$. \emph{Nearest pattern} returns the stored pattern with the largest overlap with the query, which is the limit $\beta \to \infty$ when that pattern is unique.

\subsection{Metrics and intervals}

The capacity $P_{50}$ is the load at which exact recall first falls below 50\%, interpolated linearly in $\log P$ between the two grid loads around the crossing. Intervals are 95\% percentile intervals from 2,000 cluster-bootstrap resamples: over the 50 pattern sets for the synthetic grid, over the 200 subsets for character subsets, and over the 200 draws for the whole character set. The whole-character intervals therefore measure Monte Carlo error on one fixed set of patterns, not variation across pattern sets. Because every rule sees the same queries and the classical rules use the same update orders, differences between rules are paired and can be less noisy than the individual intervals suggest.

\emph{A caution on small capacities.} When recall is about 1 at $P = 2$ and about 0 at $P = 3$, the interpolation returns $\sqrt{6} \approx 2.4$ whatever the rule. A $P_{50}$ below 3 is therefore read only as a crossing between 2 and 3 patterns, and such values are not compared with each other. Small capacities are compared through recall at $P = 3$.
